% SPDX-License-Identifier: CC-BY-4.0
\section{A second proof of the value of \texorpdfstring{\(D(n,k)\)}{D(n,k)}}\label{sec:second-proof}

This section gives the proof of the lower half of Corollary~\ref{cor:main} from
version~1 of this paper. It does not use Section~\ref{sec:lower}, and it does not
determine \(\delta(n,k,\ell)\) for individual orders. Its lower bound combines
the multiplicity Schwartz--Zippel lemma with a one-dimension reduction; the
upper bound is Theorem~\ref{thm:construction}, or the hyperplane covers for
\(k\ge2^{n-1}\). The formal proof of Lemma~\ref{lem:minimum} uses this route.

\subsection{The Schwartz--Zippel bound}

The only result used here without proof is the multiplicity Schwartz--Zippel
lemma.

\begin{theorem}[Multiplicity Schwartz--Zippel {\cite[Lemma~2.7]{DKSS13}}]\label{thm:dkss}
Let \(\F\) be a field, \(\sigma\) a nonempty finite set, \(P\in\F[x_\sigma]\) a
nonzero polynomial and \(S\subseteq\F\) a finite set. Then
\[
 \sum_{a\in S^\sigma}\mult_a(P)\le\deg P\cdot|S|^{|\sigma|-1}
\]
\lean{third-party-claims/MultiplicitySchwartzZippel.lean}.
\end{theorem}

The Lean proof follows the induction on the number of variables in~\cite{DKSS13}
(Lemma~8 of the arXiv version). It also proves the scaled form
\(|S|\sum_a\mult_a(P)\le\deg P\cdot|S|^{|\sigma|}\), which holds with no variables.

\begin{proposition}\label{prop:cover-bound}
Let \(n\ge1\) and \(k\ge0\). Every nonzero \(P\in\Ft[x_1,\ldots,x_n]\) with
\(\mult_a(P)\ge k\) for all \(a\ne\zero\) satisfies
\[
 \deg P\ge2k-\lfloor k/2^{n-1}\rfloor
\]
\lean{claims/BinaryMultiplicityCoverBound.lean}.
In particular this holds for every \(k\)-admissible \(P\).
\end{proposition}

\begin{proof}
Put \(d=\deg P\) and \(m=2^{n-1}\). The \(2^n-1\) nonzero points contribute at least
\(k(2^n-1)\) to \(\sum_{a\in\Ft^n}\mult_a(P)\), and Theorem~\ref{thm:dkss} with
\(S=\Ft\) bounds this sum by \(dm\). So \(k(2m-1)\le dm\).

Let \(q=\lfloor k/m\rfloor\), and suppose \(d+q<2k\). Then
\(dm+(q+1)m\le2km\). Since \((q+1)m>k\), this gives \(dm<2km-k=k(2m-1)\), a
contradiction.
\end{proof}

This is the deduction of the earlier note~\cite{Note}, stated without the
restriction \(k\ge2^{n-2}\) of~\cite[Question~4.3]{BBDM23}. As the note
observed, the origin condition is used only to exclude \(P=0\).

\subsection{Linear substitutions}

For a matrix \(A\in\F^{\sigma\times\sigma}\) over any field, let \(P\mapsto P\circ A\) be
the algebra map with \(x_i\mapsto\sum_jA_{ij}x_j\). The same matrix acts on
points by \((Aa)_i=\sum_jA_{ij}a_j\).

\begin{lemma}\label{lem:subst}
Let \(A,B\in\F^{\sigma\times\sigma}\) and \(P\in\F[x_\sigma]\).
\begin{enumerate}
 \item If \(P\) is homogeneous of degree \(d\), so is \(P\circ A\). Hence
 \(P\mapsto P\circ A\) commutes with taking homogeneous components, and
 \(\deg(P\circ A)\le\deg P\).
 \item \((P\circ A)\circ B=P\circ(AB)\).
 \item If \(A\) is invertible, then \(\mult_{\zero}(P\circ A)=\mult_{\zero}(P)\) and
 \(\mult_a(P\circ A)=\mult_{Aa}(P)\) for every \(a\).
\end{enumerate}
\lean{claims/OneDimensionStep.lean}
\end{lemma}

\begin{proof}
Part~1 holds because each \(x_i\) goes to a linear form. Part~2 can be checked
on the variables. For part~3, by part~2 an invertible \(A\) gives an injective
substitution. The origin order is the least \(j\) with a nonzero homogeneous
component of degree \(j\), and by part~1 this is preserved.

For the second claim, the shift commutes with substitution:
\((P\circ A)(a+z)=P(Aa+Az)=(P(Aa+z))\circ A\). So
\(\mult_a(P\circ A)=\mult_{\zero}(P(Aa+z)\circ A)=\mult_{\zero}(P(Aa+z))=\mult_{Aa}(P)\).
\end{proof}

\begin{lemma}\label{lem:linear-form}
Write \(\lambda_v=\sum_iv_ix_i\) for the linear form of a vector \(v\).
\begin{enumerate}
 \item Let \(L\in\Ft[x_\sigma]\) be nonzero with \(\deg L<2^{|\sigma|}-1\). Then
 \(\lambda_v\nmid L\) for some nonzero \(v\in\Ft^\sigma\).
 \item Over any field \(\F\), for every nonzero \(v\in\F^\sigma\) and every
 \(i_0\in\sigma\) there is an invertible \(A\) with \(\lambda_v\circ A=x_{i_0}\).
\end{enumerate}
\lean{claims/LinearFormAvoidance.lean}
\end{lemma}

\begin{proof}
(1) The ring \(\Ft[x_\sigma]\) is a unique factorization domain. Each nonzero
\(\lambda_v\) has degree one, so it is irreducible, hence prime. The only unit of \(\Ft[x_\sigma]\) is \(1\), so
distinct forms are not associated. Distinct vectors give distinct forms. If all
\(2^{|\sigma|}-1\) nonzero forms divided \(L\), their product would divide \(L\),
forcing \(\deg L\ge2^{|\sigma|}-1\).

(2) By definition \(\lambda_v\circ A=\lambda_{vA}\), with \(vA\) the row vector \(v\)
times \(A\). Extend the isomorphism between the lines \(\F v\) and \(\F e_{i_0}\)
that sends \(v\) to \(e_{i_0}\) to a linear automorphism \(g\) of \(\F^\sigma\), and
let \(A\) be the transpose of the matrix of \(g\). Then \(vA=g(v)=e_{i_0}\).
\end{proof}

\subsection{The one-dimension step}

Write the variable set as \(\sigma\cup\{x_0\}\) and a point as \((c,a')\) with
\(c\in\Ft\) and \(a'\in\Ft^\sigma\). For \(Q\in\Ft[x_0,x_\sigma]\) and \(c\in\Ft\), let
\(Q(c,x')\in\Ft[x_\sigma]\) denote the substitution \(x_0=c\).

\begin{lemma}\label{lem:restrict}
For all \(Q\), \(c\) and \(a'\), \(\mult_{a'}(Q(c,x'))\ge\mult_{(c,a')}(Q)\)
\lean{claims/OneDimensionStep.lean}.
\end{lemma}

\begin{proof}
The shift of \(Q(c,x')\) at \(a'\) is \(Q((c,a')+(w,z'))\) with \(w=0\). Its
coefficients are the coefficients of \(Q((c,a')+(w,z'))\) at monomials without
\(w\), and these have the same degree.
\end{proof}

\begin{theorem}[One-dimension step]\label{thm:step}
Let \(\sigma\) be a finite set, \(n=|\sigma|+1\) and \(k<2^n\). Let
\(P\in\Ft[x_0,x_\sigma]\) satisfy \(\mult_a(P)\ge k\) for every nonzero point \(a\)
and \(\mult_{\zero}(P)<k\). Then there is \(S\in\Ft[x_\sigma]\) with
\(\mult_{a'}(S)\ge k\) for every nonzero \(a'\), with
\(\mult_{\zero}(S)=\mult_{\zero}(P)\), and with
\[
 \deg S+1\le\deg P
\]
\lean{claims/OneDimensionStep.lean}.
\end{theorem}

In particular \(D(n-1,k)\le D(n,k)-1\) for \(n\ge2\) and \(k<2^n\). Since the
origin order is kept exactly, the same holds for the least degree with a
prescribed origin order.

\begin{proof}
Let \(\ell=\mult_{\zero}(P)<k\) and let \(L\ne0\) be the homogeneous component of
\(P\) of degree \(\ell\).

\emph{Choosing coordinates.} We have \(\deg L=\ell\le k-1<2^n-1\), so
Lemma~\ref{lem:linear-form} gives a nonzero form \(\lambda_v\nmid L\) and an
invertible \(A\) with \(\lambda_v\circ A=x_0\). Put \(Q=P\circ A\). By
Lemma~\ref{lem:subst}:
\begin{itemize}
 \item \(Q\) satisfies the hypotheses, because \(A\) permutes the nonzero points;
 \item \(\mult_{\zero}(Q)=\ell\) and \(\deg Q\le\deg P\);
 \item the degree-\(\ell\) component of \(Q\) is \(L\circ A\);
 \item \(x_0\nmid L\circ A\): otherwise substituting \(A^{-1}\) would give
 \(\lambda_v\mid L\).
\end{itemize}

\emph{The difference.} Let \(S=Q(0,x')+Q(1,x')\). Write \(Q=\sum_{t\ge0}x_0^tQ_t\)
with \(Q_t\in\Ft[x_\sigma]\). Then \(S=\sum_{t\ge1}Q_t\), since the \(t=0\) terms
cancel in characteristic two.
\begin{itemize}
 \item \emph{Degree.} \(\deg Q_t+t\le\deg Q\), so \(\deg S\le\deg Q-1\). Moreover
 \(\deg Q\ge1\), for otherwise \(S=0\), which the next item excludes.
 \item \emph{Nonzero points.} For \(a'\ne\zero\), both \((0,a')\) and \((1,a')\) are
 nonzero, so Lemma~\ref{lem:restrict} and Lemma~\ref{lem:mult-basic}(2) give
 \(\mult_{a'}(S)\ge k\).
 \item \emph{Origin order, lower bound.} Lemma~\ref{lem:restrict} gives
 \(\mult_{\zero}(Q(0,x'))\ge\ell\), and \(\mult_{\zero}(Q(1,x'))\ge k\) because
 \((1,\zero)\ne\zero\). So \(\mult_{\zero}(S)\ge\ell\).
 \item \emph{Origin order, upper bound.} The constant term of \(L\circ A\) in
 \(x_0\), namely \((L\circ A)(0,x')\), is nonzero because \(x_0\nmid L\circ A\).
 Take a monomial \(z^\beta\) with a nonzero coefficient in it. Then
 \(|\beta|=\ell\), and the same coefficient appears in \(Q(0,x')\). The
 coefficient of \(z^\beta\) in \(Q(1,x')\) vanishes, because
 \(\ell<k\le\mult_{\zero}(Q(1,x'))\). So \(z^\beta\) occurs in \(S\), and
 \(\mult_{\zero}(S)\le\ell\).
\end{itemize}
Finally \(\deg S+1\le\deg Q\le\deg P\).
\end{proof}

\subsection{The lower bound}

\begin{theorem}\label{thm:lower}
Let \(k\ge1\) and \(n\ge\lfloor\log_2k\rfloor+2\). Every \(k\)-admissible
\(P\in\Ft[x_1,\ldots,x_n]\) satisfies
\[
 \deg P\ge2k+n-2-\lfloor\log_2k\rfloor
\]
\lean{claims/BinaryMultiplicityLowerBound.lean}.
\end{theorem}

\begin{proof}
Let \(m=\lfloor\log_2k\rfloor+2\), so \(k<2^{m-1}\). Argue by induction on \(n\ge m\).

For \(n=m\), Proposition~\ref{prop:cover-bound} gives
\(\deg P\ge2k-\lfloor k/2^{m-1}\rfloor=2k\), which is the claim.

For \(n+1>m\), identify the variables with \(\{x_0\}\cup\sigma\), where \(|\sigma|=n\).
This relabelling preserves degree, and by Lemma~\ref{lem:mult-basic}(5) it
preserves multiplicities at the relabelled points. Since
\(k<2^{m-1}\le2^{n+1}\), Theorem~\ref{thm:step} gives a \(k\)-admissible \(S\) in
\(n\) variables with \(\deg S+1\le\deg P\). The induction hypothesis bounds
\(\deg S\).
\end{proof}

\subsection{The value of \texorpdfstring{\(D(n,k)\)}{D(n,k)}}

For \(k\ge2^{n-1}\), version~1 used the hyperplane covers of~\cite{BBDM23} as
the upper bound, and the formal proof still does.

\begin{theorem}[Hyperplane covers {\cite[Theorem~1.2(a) with \(d=1\), arXiv v1]{BBDM23}}]\label{thm:hyperplane}
Let \(n\ge1\) and \(k\ge2^{n-1}\), and write \(k=a2^{n-1}+b\) with
\(0\le b<2^{n-1}\). Then
\[
 P=F^a\cdot(x_1^2+x_1)^b
\]
is \(k\)-admissible and \(\deg P\le2k-\lfloor k/2^{n-1}\rfloor\)
\lean{third-party-claims/HyperplaneCoverUpperBound.lean}.
\end{theorem}

In the language of~\cite{BBDM23}, \(P\) is the product of the linear forms of
a cover by \(a\) copies of the \(2^n-1\) hyperplanes \(u\cdot x=1\), which avoid
the origin, and \(b\) copies of the parallel pair \(x_1=0\), \(x_1=1\). Their
theorem gives the least size of a cover for every \(k\ge2^{n-2}\); we use only
this upper bound for \(k\ge2^{n-1}\), in the polynomial form above, which is
what the Lean file proves. It follows from Lemma~\ref{lem:hyperplane-product}:
the multiplicity at \(a\ne\zero\) is at least \(a2^{n-1}+b=k\), the origin order
is \(b<k\), and the degree is at most \(a(2^n-1)+2b=2k-a\).

\begin{theorem}[Second proof of the value of \(D(n,k)\)]\label{thm:main-formal}
For all \(n\ge1\) and \(k\ge1\), \(D(n,k)\) takes the values of
Corollary~\ref{cor:main}. More precisely:
\begin{enumerate}
 \item for \(n\ge\lfloor\log_2k\rfloor+2\), equivalently \(k<2^{n-1}\), every
 \(k\)-admissible \(P\) has \(\deg P\ge2k+n-2-\lfloor\log_2k\rfloor\);
 \item for every \(n\ge1\), some \(k\)-admissible \(P\) has
 \(\deg P\le2k+n-2-\lfloor\log_2k\rfloor\);
 \item for \(k\ge2^{n-1}\), \(D(n,k)=2k-\lfloor k/2^{n-1}\rfloor\).
\end{enumerate}
\lean{claims/BinaryMultiplicityDegreeFormula.lean}\,\lean{claims/BinaryMultiplicityDegreeComplete.lean}
\end{theorem}

\begin{proof}
Part~1 is Theorem~\ref{thm:lower}. For \(k\ge1\), \(k<2^{n-1}\) is equivalent to
\(\lfloor\log_2k\rfloor<n-1\), which is \(n\ge\lfloor\log_2k\rfloor+2\).

For part~2, let \(J=\lfloor\log_2k\rfloor\) and \(\ell=k-2^J<k\), so that
\(k-\ell-1=2^J-1\). By Proposition~\ref{prop:menezes-polynomial} the polynomial
\(y_1^\ell g_{2^J}\) is \(k\)-admissible, with degree at most
\(n+2k-2-s_2(2^J-1)=n+2k-2-J\) by Lemma~\ref{lem:digits}(4).

For part~3, Proposition~\ref{prop:cover-bound} gives the lower bound and
Theorem~\ref{thm:hyperplane} the upper bound.
\end{proof}

